% ==============================================================================
% FFGFT: Document 246 (English)
% Title: RSG and FFGFT — A Structural Comparison
% Author: Johann Pascher
% Date: May 2026
% ==============================================================================

\section*{Abstract}

Recursive Survival Geometry (RSG, Austin 2024--2026) and Fundamental Fractal
Geometric Field Theory (FFGFT/T0, Pascher 2024--2026) address overlapping
questions — what physical structure persists, and why — from structurally
different starting points. This document compares the two frameworks along
four dimensions: what each framework takes as primitive, what it derives, where
the two accounts are complementary, and where FFGFT goes deeper in the sense of
producing numerical predictions that RSG does not. The conclusion is that the
two frameworks are asymmetrically related: FFGFT provides the scale structure
that RSG requires but does not derive; RSG provides a general selection
principle that FFGFT does not formulate explicitly. Neither is an overlay or
a bridge layer of the other.

% ------------------------------------------------------------------------------
\section{What RSG Is}
% ------------------------------------------------------------------------------

Recursive Survival Geometry is Peter Austin's framework for observable
physical structure as the differential persistence of histories under a
recursive survival-weighted filter. The central equation is:

\begin{equation}
  \frac{dS_i}{dt} = -\Gamma_i \cdot W_i \cdot S_i
\end{equation}

where $S_i$ is the survival weight of history $i$, $\Gamma_i$ is the loss
rate, and $W_i$ is the survival weight. What we observe as matter, radiation,
or structure are the histories that persist under this filter. Histories that
fail to persist below the threshold are not observed.

RSG's ontological primitive is the \emph{history} — a path through
configuration space — and the \emph{survival filter} — the operator that
weights persistence. The filter is general: RSG does not commit to a specific
Lagrangian, a specific field content, or a specific set of coupling constants.
It describes the structural logic of selection without specifying what is being
selected from.

\subsection{What RSG derives}

RSG derives the structural conditions under which a history can persist: it
must survive repeated application of the filter, it must maintain closure
under the recursive process, and it must accumulate weight $W_i$ above
threshold. The formalism produces a general account of why some configurations
persist and others do not.

\subsection{What RSG does not derive}

RSG does not derive the values of physical constants. The loss rates $\Gamma_i$,
the survival weights $W_i$, and the threshold conditions are framework
parameters. RSG does not explain why the electron mass is $0.511\,\text{MeV}$,
why the fine structure constant is $\alpha \approx 1/137$, or why the CMB
temperature is $2.725\,\text{K}$. These are inputs to RSG's filter, not
outputs of it.

% ------------------------------------------------------------------------------
\section{What FFGFT Is}
% ------------------------------------------------------------------------------

Fundamental Fractal Geometric Field Theory (T0/Time-Mass Duality) takes the
4-dimensional torus $T^4 = (\mathbb{R}/2\pi\mathbb{Z})^4$ with winding
structure $(n_\theta, n_\phi, n_3, n_4)$ as its ontological primitive. The
foundational relation is $\tilde{T} \cdot m = 1$ (time-mass duality). A
single dimensionless parameter $\xi = 4/30000$ is determined from three
measured mass ratios via the unique winding exponent $\kappa = 7$.

The T0 Lagrangian is:

\begin{align}
  \mathcal{L}_{\text{T0}} &= -\tfrac{1}{4}F_{\mu\nu}F^{\mu\nu}
    + \bar{\psi}_\ell(i\gamma^\mu D_\mu - m_\ell)\psi_\ell \nonumber\\
    &\quad + \tfrac{1}{2}(\partial_\mu \Delta m)(\partial^\mu \Delta m)
    - \tfrac{1}{2}m_T^2 \Delta m^2
    + \xi \cdot m_\ell \cdot \bar{\psi}_\ell \psi_\ell \cdot \Delta m
\end{align}

where $\Delta m(x,t)$ is the dynamical mass field and $\xi \cdot m_\ell$ is
the vertex coupling.

\subsection{What FFGFT derives}

From $\xi = 4/30000$ alone:

\begin{itemize}
  \item Lepton and quark mass ratios via $\kappa=7$ recursion (error $<1\%$, no fitting)
  \item Fine structure constant $\alpha$ (Doc~009/137)
  \item CMB temperature $T_\text{CMB}$ (deviation $\Delta = 0.0002\%$, Doc~025)
  \item Cosmological constant $\Lambda = 1.34\times10^{-122}$ (Planck units, Doc~182)
  \item Hubble constant $H_0 = 66.2\,\text{km/s/Mpc}$ (Doc~182)
  \item Geometric path-lengthening redshift mechanism (FEM simulation, Doc~026)
\end{itemize}

\subsection{What FFGFT does not formulate}

FFGFT does not have an explicit \emph{selection principle} of the RSG type. It
derives the scale structure of what exists, but it does not formulate why
those structures persist as observable histories rather than being absorbed
into the vacuum. The persistence question — why lepton-type histories survive
and not others — is answered implicitly through the winding structure of $T^4$,
but not through a dedicated survival-filter formalism.

% ------------------------------------------------------------------------------
\section{The Structural Comparison}
% ------------------------------------------------------------------------------

\renewcommand{\arraystretch}{1.5}
{\small
\begin{longtable}{p{3.5cm}p{5.0cm}p{5.0cm}}
\toprule
\textbf{Dimension} & \textbf{RSG (Austin)} & \textbf{FFGFT (Pascher)} \\
\midrule
\endfirsthead
\toprule
\textbf{Dimension} & \textbf{RSG (Austin)} & \textbf{FFGFT (Pascher)} \\
\midrule
\endhead
\midrule\multicolumn{3}{r}{\footnotesize\itshape continued\ldots}\\\endfoot
\bottomrule\endlastfoot

\textbf{RA4 Primitive} &
History + survival filter. No committed field content or Lagrangian. &
$T^4$ winding geometry. Committed Lagrangian: $\mathcal{L}_{\text{T0}}$. \\

\midrule

\textbf{RA3 Admissibility} &
Histories must survive recursive application of the filter above threshold.
General condition, no fixed coupling constants. &
Only winding configurations consistent with $\xi$ are admissible.
$D_f = 3-\xi$, Z3$^3$ constraint fixes three generations. \\

\midrule

\textbf{RA2 Representation} &
Weighted history space with survival measure $S_i$. No fixed Hilbert space. &
$L^2(T^4)$; Fourier modes with winding numbers as natural basis;
QM operator algebra derived as projection (Doc~230/231). \\

\midrule

\textbf{RA1.5 Operationalization} &
$\Gamma_i$, $W_i$ as framework parameters — not derived from a fixed input.
Scale of the filter undefined without additional input. &
All SM constants from $\xi$ alone: masses, $\alpha$, $T_\text{CMB}$,
$\Lambda$, $H_0$. No fitting. \\

\midrule

\textbf{RA1 Predictions} &
Structural persistence conditions. No specific numerical predictions
for particle masses or cosmological constants. &
Concrete numerical values for all SM constants and cosmological
parameters. Explicit falsification criteria. \\

\midrule

\textbf{Redshift} &
Not addressed explicitly in RSG core. &
Geometric path-lengthening in $\xi$-field, FEM simulation confirms
frequency independence. Indistinguishable from metric expansion
from inside (Doc~026). \\

\midrule

\textbf{Selection principle} &
Explicit: survival filter $\Gamma_i W_i$ defines which histories persist. &
Implicit: winding number admissibility. No dedicated persistence
operator formulated separately from the Lagrangian. \\

\midrule

\textbf{Scale derivation} &
Not present: $\Gamma_i$, $W_i$ require external input for specific values. &
Complete: $\xi$ from three mass ratios; all scales follow without
additional input. \\

\end{longtable}
}

% ------------------------------------------------------------------------------
\section{Is RSG Explainable by FFGFT?}
% ------------------------------------------------------------------------------

Partially. FFGFT provides the scale structure that RSG requires as input.
If $\Gamma_i = \xi \cdot m_i$ (coupling proportional to mass, as in the
$\mathcal{L}_{\text{T0}}$ vertex), then RSG's filter parameters are
determined by FFGFT's coupling structure. In that reading, RSG's survival
filter is the phenomenological description of what FFGFT's Lagrangian
produces at the level of observable history selection.

What FFGFT does not supply is RSG's general formulation of the selection
principle as a framework-independent operator. RSG's survival filter applies
to any history regardless of whether the underlying dynamics comes from FFGFT,
from standard QFT, or from some other field theory. That generality is
RSG's contribution; FFGFT's specific coupling structure is one instantiation
of it.

% ------------------------------------------------------------------------------
\section{Is FFGFT Explainable by RSG?}
% ------------------------------------------------------------------------------

No. RSG does not derive numerical values for particle masses, coupling
constants, or cosmological parameters. The electron mass $m_e = 0.511\,\text{MeV}$,
the ratio $m_p/m_e = 1836.15$, the fine structure constant $\alpha \approx
1/137$ — none of these follow from RSG's survival filter without specifying
$\Gamma_i$ and $W_i$ externally. RSG's formalism is consistent with any
value of these constants; it does not predict them.

FFGFT, by contrast, fixes $\xi$ from three measured ratios and derives
all others. This is a stronger claim, not a weaker one. FFGFT is not a
computational overlay of RSG — it is a competing and more constrained
answer to the question of what the physical constants are and why.

% ------------------------------------------------------------------------------
\section{Which Framework Is More Complete?}
% ------------------------------------------------------------------------------

The question of completeness depends on what ``complete'' means.

If completeness means \emph{numerical predictive power} — the ability to
derive specific values for observables from minimal input — then FFGFT is
more complete. It produces quantitative predictions that RSG cannot make
without external specification of its parameters.

If completeness means \emph{generality of the selection principle} — the
ability to characterise which structures persist across different underlying
dynamics — then RSG is broader. Its survival filter applies regardless of
the specific field theory; FFGFT's winding structure is one particular
realisation of a persistent configuration, not a general account of persistence.

The two frameworks are therefore asymmetrically complementary:

\begin{itemize}
  \item FFGFT provides the \emph{scale structure} that RSG's filter requires
        as input but does not derive.
  \item RSG provides the \emph{selection-principle language} that FFGFT
        instantiates but does not formulate as a general operator.
\end{itemize}

Neither framework is an overlay or bridge layer of the other. They are
two distinct frameworks that address adjacent but non-identical questions,
with partial structural overlap at the level of persistence conditions.

% ------------------------------------------------------------------------------
\section{The Asymmetry in One Sentence}
% ------------------------------------------------------------------------------

\begin{keyresult}[Structural Asymmetry]
RSG asks: \emph{which histories survive?} — without specifying the scale.
FFGFT asks: \emph{what are the scales?} — without formulating a general
persistence operator. RSG needs FFGFT's scales to be quantitative; FFGFT
needs RSG's language to explain why winding configurations are selected.
Neither reduces to the other.
\end{keyresult}

% ------------------------------------------------------------------------------
\section{On Visualisation: Two Parallel Entries}
% ------------------------------------------------------------------------------

The structural analysis in this document yields a clear requirement for any
comparative visualisation of the two frameworks.

FFGFT is not a computational overlay of RSG and not a bridge layer within
RSG's architecture. It is a standalone framework with its own RA4 primitive
($T^4$ torus), its own Lagrangian ($\mathcal{L}_{\text{T0}}$), its own
parameter ($\xi$), and its own independent numerical predictions that RSG
does not supply.

A structurally correct visualisation therefore shows:

\begin{itemize}
  \item \textbf{Two parallel entries at the same structural level} —
        RSG and FFGFT as independent frameworks side by side, not in
        a hierarchy.
  \item \textbf{A $\xi$-path as its own track} — the $\xi$ parameter
        and the scales derived from it form an independent line in
        the visualisation, separate from RSG's filter parameters
        $\Gamma_i$ and $W_i$.
  \item \textbf{Correspondence notes at the contact points} — where
        FFGFT's $\xi \cdot m_i$ determines the value of RSG's $\Gamma_i$,
        and where RSG's persistence language names what FFGFT's
        winding structure implicitly selects.
\end{itemize}

This architecture — two co-equal entries with explicit correspondence
arrows — is the only representation that does justice to both frameworks:
RSG's generality of the selection principle and FFGFT's completeness of
scale derivation remain both visible, without either being subordinated
to the other.

% ------------------------------------------------------------------------------
\section{Symbol List}
% ------------------------------------------------------------------------------

{\small
\begin{tabular}{p{3cm}p{10.5cm}}
\toprule
\textbf{Symbol} & \textbf{Meaning} \\
\midrule
$S_i$ & RSG survival weight of history $i$ \\
$\Gamma_i$ & RSG loss rate for history $i$ \\
$W_i$ & RSG survival weight factor \\
$\xi = 4/30000$ & FFGFT dimensionless coupling constant; fixed by $\kappa=7$ \\
$T^4$ & 4-dimensional torus; FFGFT ontological primitive \\
$\kappa = 7$ & Unique winding exponent from e-p-$\mu$ mass triangle \\
$\mathcal{L}_{\text{T0}}$ & T0 Lagrangian (5 terms; Doc~180) \\
$\Delta m(x,t)$ & Dynamical mass field in T0 \\
$D_f = 3-\xi$ & Fractal dimension of physical space ($\approx 2.99987$) \\
$L^2(T^4)$ & FFGFT Hilbert space \\
$\alpha$ & Fine structure constant; derived from $\xi$ (Doc~009) \\
$\Lambda$ & Cosmological constant from $\xi$: $1.34\times10^{-122}$ (Doc~182) \\
$H_0$ & Hubble constant from $\xi$: $66.2\,\text{km/s/Mpc}$ (Doc~182) \\
RA4--RA1 & Layers of Representational Architecture 2.1\\
          & \quad (Guevara Calderón 2026) \\
RSG & Recursive Survival Geometry (Austin 2024--2026) \\
FFGFT & Fundamental Fractal Geometric Field Theory (Pascher 2024--2026) \\
\bottomrule
\end{tabular}
}

\vspace{1em}
\noindent\textbf{Key references:}
Doc~009 ($\xi$, $\kappa=7$) ·
Doc~025 (CMB) ·
Doc~026 (redshift, FEM) ·
Doc~180 (Lagrangian, $\kappa=7$) ·
Doc~182 ($\Lambda$, $H_0$) ·
Doc~230/231 (Hilbert space bridge) ·
Doc~245 (FFGFT RA diagnostic) ·
\href{https://github.com/jpascher/T0-Time-Mass-Duality}{GitHub} ·
\href{https://doi.org/10.5281/zenodo.20117635}{Zenodo v1.1.0}
